This study provides evidence for a persistent asymmetry in early diffusion modeling. Across FibVID and the Twitter15/16 benchmark, early diffusion signals are substantially more informative for reach prediction than for veracity. On FibVID, the main harmonized $k$-window analysis yields strong reach prediction but only modest veracity discrimination; the Twitter benchmark shows the same broad pattern, though with somewhat stronger veracity performance. The most defensible reading is therefore not that early diffusion can determine truth reliably, but that it can support weak misinformation-risk screening and much stronger reach prioritization.

Methodologically, the results also suggest that the structure of the observation problem matters. Size-based windows remain the stronger primary design, time-window analyses are useful but weaker robustness checks, and the native graph representation recovers somewhat more veracity signal without replacing the main cascade-compatible specification. In that sense, the paper's contribution is not only a set of performance numbers, but also a clearer account of which parts of early diffusion are most informative for which task.

These conclusions remain bounded. The estimates are predictive rather than causal, conditional on explicit window definitions, and partly dependent on how propagation structure is represented. Within those bounds, however, the empirical message is clear: early diffusion is much better at indicating \emph{how far} a cascade may go than \emph{whether} it is false. For early-warning systems, this implies a staged use of diffusion signals: cautious veracity screening early on, and stronger prioritization of likely high-impact cascades as more structure becomes visible.
